%\addcontentsline{toc}{chapter}{Abstract}

\begin{abstract}
% Text of the Abstract. No more than 300 words.
% pT
% Bjorken x 7%-9%
% x ratio less than 5%
\vspace{15pt}

One of the fundamental questions in physics is the matter-antimatter asymmetry of the Universe, which is being explored through charge-parity violation searches in neutrino oscillation experiments. The next-generation oscillation experiments, such as DUNE, entering the intensity and precision frontiers, need to constrain the systematic uncertainties arising from our understanding of (anti)neutrino-nucleus scattering to within a few percent to achieve their physics goals. MINERvA, a dedicated (anti)neutrino-nucleus scattering experiment, plays a crucial role in constraining these uncertainties. This thesis presents MINERvA's first high-statistics direct measurement of nuclear dependence in charged-current inclusive antineutrino scattering on carbon, hydrocarbon, iron, and lead as a function of antimuon transverse momentum $p_{\text{T}}$ and Bjorken $x$.

The analysis utilises interactions with a mean antineutrino energy of approximately $\sim$6~GeV, which have a reconstructed $\mu^{+}$ with a scattering angle less than 17$^\circ$ relative to the antineutrino beam and an antimuon energy of 2--20 GeV. The measured per-nucleon differential cross-sections are reported with a precision of 7--9\%, while the cross-section ratios of carbon, iron, and lead to hydrocarbon have uncertainties of about 5\% or less. The measured cross-sections for iron and lead indicate strong suppression at low $p_{\text{T}}$ and Bjorken $x$, along with an enhancement of the cross-section at high $p_{\text{T}}$. These effects are observed to be more pronounced with the increasing size of the target nucleus and are not reproduced by the underlying simulation prediction. Comparisons to alternative models used in current (anti)neutrino interaction generators are included, showing some improvements in the modelling over the base prediction model, yet they are still unable to fully reproduce the observed nuclear dependence in this analysis.

Importantly, the analysis provides a direct test of nuclear effects in inclusive antineutrino scattering, with the main contributions coming from resonant pion production, deep inelastic scattering, and the transition region between these two interaction channels, which will be the major contributors in DUNE. This measurement also represents one of the largest antineutrino datasets in this energy regime analysed to date. 

\end{abstract}